Consumer headsets and smartwatches offer a route to touch-free laptop interaction, but a useful interface requires more than high classification accuracy on isolated gestures. We describe a spatial volume-control system in which head orientation selects a target, separate smartwatch pinch-start and pinch-release detections delimit a grab, and wrist rotation adjusts a bounded output. Photoplethysmography and inertial signals are processed on the laptop, while explicit cancellation and mode policies constrain actuation. We present a source-audited account that separates historical camera-labelled experiments from a later manually revised development cohort. The latter directory contains 37 exports but only 28 distinct complete PPG--IMU stream pairs; 14 exports contain explicit optical-sample annotation overrides. Replaying a stored pinch-in model over its eight listed evaluation recordings with the trainer's own detection code reproduced its stored metrics exactly: 47 of 57 annotated starts caught, no false activations over 0.48 h of eligible no-gesture time (95\% upper bound 6.2 per hour) and a median latency of 442 ms, although the everyday-activity session came within 0.02 of the threshold. Historical camera-label evaluation instead caught 161 of 372 starts in the selected leave-one-day-out configuration, revealing marked day variation. These results establish a reproducible prototype and expose the importance of annotation provenance, source-group isolation and causal event accounting; they do not establish population-level or all-day reliability. We specify the prospective, matched-modality and interaction evaluations required to test those stronger claims.
